\section{Introduction}
\label{sec:introduction}

The web was conceived as a decentralized hypertext system, but the practical web of the past two decades has consolidated around a small number of centralized service providers. Communication, computation, storage, and identity flow through their infrastructure, and as a consequence so does censorship, surveillance, and de-platforming risk. The technical question this paper addresses is how to recover the application-level functionality of those services --- group chat, real-time collaboration, social feeds, reputation, authentication --- without recreating the central choke points.

Two broad families of decentralized systems have been deployed at scale. Content-addressed storage networks such as the original Freenet \citep{clarke2001freenet} and IPFS \citep{benet2014ipfs} provide name-by-hash access to immutable data, but they offer no native mechanism for evolving shared state, so applications must layer their own consistency on top. Blockchains \citep{nakamoto2008bitcoin} provide programmable, mutable state, but the global linearizable consensus they offer is paid for in latency, throughput, and replication cost; building responsive interactive applications on top of them is awkward.

This paper describes the architecture of the new Freenet: a peer-to-peer platform whose central abstraction is a programmable, eventually-consistent key-value cell --- the \emph{contract} --- replicated across a small-world peer-to-peer network. A contract is a WebAssembly module that defines, for its application, a commutative-monoid merge function on state and an efficient summary/delta synchronization protocol. The platform provides routing, replication, and a uniform consistency model; the application provides the algebra. Alongside contracts, Freenet supports \emph{delegates}: actor-style WebAssembly agents that hold private state on a user's device and communicate by attributed messages. Public contract state and private delegate state are explicitly separated.

\subsection{Contributions}

The paper makes the following architectural contributions, each described in its own section:

\begin{enumerate}[leftmargin=1.4em, itemsep=2pt]
  \item \textbf{Contract-defined commutative monoids on application state} (\cref{sec:monoids}). Conventional Conflict-free Replicated Data Types (CRDTs) \citep{shapiro2011crdt} expose a fixed catalogue of lattices to applications. Freenet generalizes this by letting each contract carry its own merge function, validity predicate, and identity element in sandboxed WebAssembly. The runtime is generic over the algebra; correctness follows from contract-level properties rather than runtime-enforced ones.

  \item \textbf{Summary/delta synchronization parameterized by the contract} (\cref{sec:sync}). Each contract exposes a triple $(\summarize, \getdelta, \applydelta)$ that lets two peers reconcile their replicas using a compact, application-defined summary instead of an exchange of full state. The mechanism subsumes both Merkle anti-entropy and operation-based update propagation, without requiring a causal-broadcast layer.

  \item \textbf{Adaptive routing via isotonic regression} (\cref{sec:adaptive}). Freenet's request router does not forward greedily to the topologically closest peer. It maintains, per neighbor, a monotone non-parametric estimator of latency, failure probability, and throughput as a function of distance-to-target, and forwards each request to the neighbor with the lowest expected completion cost. The result is performance-aware small-world routing that adapts to heterogeneous peer capacity.

  \item \textbf{Accept-only-at-terminus topology construction} (\cref{sec:smallworld}). The small-world property is not engineered by inserting long links at chosen lengths in the style of Kleinberg \citep{kleinberg2000small}; instead it emerges from a single connection-acceptance rule: a join request is forwarded greedily until it reaches a peer with no neighbor closer to the target than itself, and only that peer accepts the connection. Local clustering and the long-range links needed for logarithmic-diameter routing both arise from this rule applied to bootstrap and to ordinary request traffic.

  \item \textbf{The delegate model for private state and sensitive operations} (\cref{sec:delegates}). Delegates are sandboxed WebAssembly agents, local to a user's device, that hold secrets and respond to authenticated messages from contracts, user interfaces, or other delegates. They give the platform a built-in trust boundary between applications that handle private keys and the user interfaces that drive them, with cross-device replication via shared-secret contracts.
\end{enumerate}

The combination of these mechanisms supports a class of decentralized applications --- group chat \citep{river}, collaborative documents, reputation systems --- whose consistency requirements are accommodated by per-application merge semantics rather than by network-wide linearizability. We position the system relative to prior work in \cref{sec:related}, discuss the consistency boundary it occupies and its limits in \cref{sec:discussion}, and conclude in \cref{sec:conclusion}.
